\section{The part tree \& mass properties}\label{sec:parts}

Everything the simulator knows about a rocket's mass --- how heavy it is, where its center of
mass (CM) sits, how hard it resists rotation --- comes from one data structure: a composite
tree of \code{Part} nodes under \srcf{model/parts/}. Each node is simultaneously one physical
component and the root of a sub-tree, so it carries both its own mass properties and the lazily
cached composite of itself plus every descendant \src{model/parts/Part.h}{26}. This chapter
walks the tree from the outside in: ownership and identity (\cref{sec:parts-tree}), the five
concrete part types (\cref{sec:parts-catalog}), the two-pass mass-property algorithm at the
chapter's core (\cref{sec:parts-massprops}), the caching that makes those properties
time-aware (\cref{sec:parts-caching}), the closed-form tensor library
(\cref{sec:parts-tensors}), and the diagnostics gate on the composite readers
(\cref{sec:parts-gating}). How a child's position is \emph{stored and resolved} --- the
\code{StationLink}, the resolver, and the overlap sweep --- is the subject of
\cref{sec:placement}; this chapter treats placement as an opaque intent value riding alongside
each child. The motor's internals live in \cref{sec:motors}. All quantities are SI --- meters,
kilograms --- with one deliberate exception flagged in \cref{sec:parts-catalog}.

\subsection{The composite tree}\label{sec:parts-tree}

\code{Part} is the abstract node of a classic Composite pattern
(\cref{fig:parts-uml-tree}). A node owns its children as a vector of
(\code{shared\_ptr<Part>}, \code{StationLink}) pairs in attachment order
\src{model/parts/Part.h}{288}: the smart pointer carries ownership of the child sub-tree, and
the \code{StationLink} beside it stores the child's \emph{placement intent} --- a
feature-to-feature relationship, never a resolved coordinate
(\cref{sec:placement}). Each node also keeps a non-owning back-pointer to its parent, used to
propagate cache-invalidation flags up the tree
\src{model/parts/Part.h}{222}. Two members are pure virtual, and they define what a concrete
part must supply: \code{typeName()}, the stable string tag that doubles as the factory key and
the design-file type attribute \src{model/parts/Part.h}{181}, and \code{cloneShallow()}, the
type-preservation hook for deep copies \src{model/parts/Part.h}{217}.

\tikzfig{uml-part-tree}{The \code{Part} composite tree. The abstract base owns children as
(\code{shared\_ptr}, \code{StationLink}) pairs in attachment order
\srccap{model/parts/Part.h}{288} and keeps a non-owning parent back-pointer for upward cache
invalidation \srccap{model/parts/Part.h}{222}. Five concrete types override the pure virtuals
\code{typeName()} and \code{cloneShallow()}; \code{PartFactory} is the sanctioned construction
path for the four geometry types, while \code{Motor} is attached only through
\code{RocketModel::setMotorModel()} \srccap{model/parts/PartFactory.cpp}{73}.}{fig:parts-uml-tree}

Identity is the id, never the name. Every part gets a process-unique
\code{Part::Id} (a \code{std::uint64\_t}) from a relaxed atomic counter that starts at 1,
reserving 0 as ``none'' \src{model/parts/Part.cpp}{26}; the id is assigned fresh in every
constructor --- including the copy constructor --- and never reassigned
\src{model/parts/Part.h}{224}. Names are human-facing labels and need not be unique
\src{model/parts/Part.h}{176}; ids are deliberately not serialized, because they are only
meaningful within a single run \src{model/parts/Part.h}{171}. The placement diagnostics of
\cref{sec:placement} identify parts by id for exactly this reason: names collide and pointers
dangle.

\begin{designrule}[A part lives at one place in one tree]
Copy and move assignment are deleted, and the protected copy constructor suppresses the
implicit move, so a \code{Part} cannot be copied, moved, or sliced at the value level
\src{model/parts/Part.h}{66}. A part enters a tree only by moving its owning
\code{shared\_ptr} into \code{addChildPart()}, and is duplicated only via \code{clone()} ---
``the only way to duplicate a part'' \src{model/parts/Part.h}{193}. Aliasing a node into two
trees, or double-attaching it within one, is unrepresentable rather than merely discouraged.
\end{designrule}

\code{addChildPart(child, link)} enforces that rule with three guards, each a logged no-op
rather than an exception: a null child \src{model/parts/Part.cpp}{73}; a child that already
has a parent, with the log message pointing at the remedy (``clone() it or detach first'')
\src{model/parts/Part.cpp}{78}; and a child that is this part or one of its ancestors, which
would form both a graph cycle and an ownership cycle through the child \code{shared\_ptr}s
\src{model/parts/Part.cpp}{84}. Adoption is then re-parenting followed by a move --- no copy,
so the child's dynamic type and id are preserved
\src{model/parts/Part.cpp}{97}. Detachment is the inverse: \code{removeChildById()} moves the
owning pointer out, clears the child's parent, and returns the intact sub-tree to the caller
\src{model/parts/Part.cpp}{277}. Both operations end by invalidating \emph{both} caches ---
composite mass properties and resolved placement --- up the whole ancestor chain
\src{model/parts/Part.cpp}{100}; removal drives the rebuild through the dirty flag rather than
the mass-delta gate of \cref{sec:parts-caching}, so even detaching a zero-mass sub-tree
refreshes the composite \src{model/parts/Part.h}{205}.

Duplication is a two-level protocol (\cref{lst:parts-clone}). \code{cloneShallow()} copies
just one node with the correct dynamic type; each concrete class implements it as
\cppinline|new X(*this)| through a protected, defaulted copy constructor (e.g.\
\src{model/parts/BodyTube.h}{63}), and the base copy constructor copies the node's own mass
properties, mints a fresh id, and leaves the copy parentless and childless
\src{model/parts/Part.cpp}{55}. \code{clone()} then recurses polymorphically over the
children, copying each \code{StationLink} verbatim --- the placement intent is plain data and
clones exactly. The result is a deep, independent, type-preserving copy of the whole sub-tree,
pinned per concrete type by the test suite \src{model/tests/PartTests.cpp}{252}.

\begin{listing}[htbp]
\begin{minted}{cpp}
std::shared_ptr<Part> Part::clone() const
{
   std::shared_ptr<Part> copy = cloneShallow(); // this node: correct dynamic type, fresh id, no kids
   for(const auto& [child, link] : childParts)
   {
      std::shared_ptr<Part> childCopy = child->clone(); // recurse polymorphically (no slicing)
      childCopy->parent = copy.get();
      copy->childParts.emplace_back(std::move(childCopy), link); // the placement intent clones verbatim
   }
   return copy;
}
\end{minted}
\caption{Deep copy without slicing \srccap{model/parts/Part.cpp}{106}: the pure-virtual
\code{cloneShallow()} reproduces each node's dynamic type, the recursion reproduces the
structure, and the copied \code{StationLink}s reproduce the placement intent. Every cloned
node receives a fresh id.}\label{lst:parts-clone}
\end{listing}

After construction a part is close to immutable: the only setters are \code{setMass()} and
\code{setI()}, both of which flag the ancestor chain for recompute
\src{model/parts/Part.h}{74}; there are no geometry setters on any concrete type ---
changing a dimension means building a replacement part (and losing the id). Geometry is
exposed read-only through a family of virtuals: \code{getLength()} (a part occupies
$z \in [-L, 0]$ in its own \zfwd{} frame, fore plane at the local origin)
\src{model/parts/Part.h}{130}, the silhouette profiles
\code{radiusOuterAt()}/\code{radiusInnerAt()} \src{model/parts/Part.h}{135},
\code{axialLength()}, \code{isSolid()}, and \code{stationAt()} \src{model/parts/Part.h}{144},
the solid-host capacity rule \code{innerCapacityAt()} \src{model/parts/Part.h}{158}, and the
frontal-disc reference areas \src{model/parts/Part.h}{162}. \Cref{sec:placement} gives these
their semantics; the catalog below records what each concrete type returns.

The sanctioned construction path is \code{PartFactory}. One type-agnostic \code{PartParams}
struct spans every factory-buildable part, every geometry field a \code{std::optional} so a
missing key differs from a supplied zero \src{model/parts/PartFactory.h}{27}; the CLI's
key=value parser, \code{makePart()}, and the reflection function \code{params()} all share
it, ``so the field vocabulary can't drift between the front-ends and the on-disk format''
\src{model/parts/PartFactory.h}{23}. \code{makePart(type, p)} dispatches on the exact
\code{typeName()} string \src{model/parts/PartFactory.cpp}{27}: a missing required field
throws a type-named \code{std::invalid\_argument} \src{model/parts/PartFactory.cpp}{16},
constructor range validation propagates, and an unknown type --- including \code{"Motor"},
attached through \code{RocketModel::setMotorModel()}, never built directly --- throws as well
\src{model/parts/PartFactory.cpp}{73}. The factory always passes a zero CM offset:
``placement is the caller's attach offset, not a baked-in CM''
\src{model/parts/PartFactory.h}{54}. \code{params()} reflects a part's own scalars back
through a \code{dynamic\_cast} chain \src{model/parts/PartFactory.cpp}{79}, and the round
trip reproduces mass and full composite CM/inertia to within four ulps
(\code{EXPECT\_DOUBLE\_EQ}) \src{model/tests/PartFactoryTests.cpp}{139}. New part headers register in one place, the
umbrella header \srcf{model/parts/Parts.h} \src{model/parts/Parts.h}{4}.

\subsection{The concrete catalog}\label{sec:parts-catalog}

Five concrete types derive from \code{Part} (\cref{tab:parts-catalog}). The four geometry
types compute their mass, per-unit-mass inertia tensor, and CM offset from geometry plus
density at construction, delegating the tensor algebra to \code{model::InertiaTensors}
(\cref{sec:parts-tensors}); their constructors validate ranges and throw
\code{std::invalid\_argument} on non-physical geometry, so a constructed part is always
internally consistent. The \code{Motor} takes its mass and dimensions from the wrapped
catalog entry instead.

\begin{table}[htbp]
\centering\small
\begin{tabularx}{\linewidth}{@{}l>{\raggedright\arraybackslash}X>{\raggedright\arraybackslash}Xll@{}}
\toprule
\code{typeName()} & Geometry parameters & Own mass & Tensor & $C_{N\alpha}$ \\
\midrule
\code{BodyTube} & $r_i$, $r_o$, $L$, $\rho$ & $\rho\,\pi(r_o^2 - r_i^2)L$ & \code{Tube} & $0$ \\
\addlinespace
\code{NoseCone} & $R$, $L$, wall $t$, $\rho$, solid flag &
  $\rho\,\tfrac{1}{3}\pi R^2 L$ (solid); $\rho\,\pi R\sqrt{R^2{+}L^2}\,t$ (shell) &
  \code{SolidCone} / \code{ConicalShell} & $2$ \\
\addlinespace
\code{FinSet} & $N$, $c_r$, $c_t$, $s$, $X_t$, $t$, $r_b$, $\rho$ &
  $N\rho\,\tfrac{1}{2}(c_r{+}c_t)\,s\,t$ & \code{TrapezoidalFinSet} & fin term \\
\addlinespace
\code{HollowSphere} & $r_i$, $r_o$, $\rho$ & $\rho\,\tfrac{4}{3}\pi(r_o^3 - r_i^3)$ &
  \code{HollowSphere} & inert \\
\addlinespace
\code{Motor} & wraps a \code{MotorModel} & \code{mm.getMass($t$)} &
  cylinder from motor $d$, $L$ & inert \\
\bottomrule
\end{tabularx}
\caption{The concrete part catalog at the pinned commit \code{254cd41}. Mass closed forms:
\srccap{model/parts/BodyTube.cpp}{28}, \srccap{model/parts/ConicalNoseCone.cpp}{32},
\srccap{model/parts/FinSet.cpp}{46}, \srccap{model/parts/HollowSphere.cpp}{42},
\srccap{model/parts/Motor.h}{40}. The $C_{N\alpha}$ column is each part's Barrowman
normal-force-slope contribution, analyzed in \cref{sec:aero}.}\label{tab:parts-catalog}
\end{table}

\begin{description}
\item[\code{BodyTube}] --- a uniform-density hollow circular cylinder, axis on $z$; a solid
  rod when $r_i = 0$. The constructor validates $0 \le r_i < r_o$, $L > 0$, $\rho > 0$
  \src{model/parts/BodyTube.cpp}{22}; the CM sits at the geometric center
  \src{model/parts/BodyTube.h}{25} and the silhouette
  is constant walls \src{model/parts/BodyTube.h}{55}. It contributes a reference area
  $\pi r_o^2$ \src{model/parts/BodyTube.h}{52} and a wetted area $2\pi r_o L$ held ready for
  a future skin-friction drag term \src{model/parts/BodyTube.h}{51}; aerodynamically it
  reports $C_{N\alpha} = 0$ --- a constant-diameter body produces no normal force, so it
  drops out of the center-of-pressure average \src{model/parts/BodyTube.cpp}{38}.
\item[\code{ConicalNoseCone}] --- a straight right circular cone, either solid (the default)
  or a thin lateral shell (open base, wall thickness $t \ll R$)
  \src{model/parts/ConicalNoseCone.h}{22}. The two variants differ in everything derived:
  mass \src{model/parts/ConicalNoseCone.cpp}{32}, tensor
  (\code{SolidCone} vs.\ \code{ConicalShell} \src{model/parts/ConicalNoseCone.cpp}{44}), and
  CM --- the solid centroid is $L/4$ from the base, the shell's $L/3$, reported as offsets
  $-L/4$ and $-L/6$ from the component middle in the \zfwd{} frame
  \src{model/parts/ConicalNoseCone.cpp}{49}. The silhouette is the linear taper
  $R\cdot(-z/L)$ with a zero-length divide guard \src{model/parts/ConicalNoseCone.h}{61}, and
  \code{isSolid()} returns the \emph{authored} flag, because a shell cone's bore vanishes at
  the tip and bore inference would misclassify it
  \src{model/parts/ConicalNoseCone.h}{55}. Its Barrowman contribution is the classic nose
  $C_{N\alpha} = 2$ at its own base area, rescaled to the shared reference area, with the
  center of pressure at $2L/3$ aft of the tip \src{model/parts/ConicalNoseCone.cpp}{57}. It is
  the only nose profile in the tree --- no ogive, parabolic, or Haack shapes exist at this
  commit (\cref{sec:incomplete}).
\item[\code{FinSet}] --- $N \ge 3$ identical symmetric trapezoidal flat-plate fins, modeled
  as \emph{one analytic leaf}, not $N$ child parts \src{model/parts/FinSet.h}{22}. The
  geometry vocabulary is the standard planform set: root chord $c_r$, tip chord $c_t$ (zero
  allowed --- a triangular fin), semi-span $s$, plate thickness, mounting body radius $r_b$,
  and the leading-edge (LE) sweep \emph{length} $X_t$ in meters --- an axial offset, not an
  angle \src{model/parts/FinSet.h}{45}. Its axial extent is the root chord
  \src{model/parts/FinSet.h}{61}, and its CM lies on the axis at the trapezoid's axial mass
  centroid, reported as $x_c - c_r/2$ from the middle \src{model/parts/FinSet.cpp}{54}. The
  envelope is deliberately honest: \code{radiusOuterAt()} reports only the body disc $r_b$,
  never $r_b + s$ \src{model/parts/FinSet.h}{73}, and the reference area is the body disc
  $\pi r_b^2$ --- fins do not inflate the reference area
  \src{model/parts/FinSet.h}{71}. A fin count below three is genuinely anisotropic and
  unsupported: the constructor warns and falls back to the $N \ge 3$ isotropic tensor rather
  than throwing \src{model/parts/FinSet.cpp}{36}. Its Barrowman contribution is the standard
  $N$-fin term with the mid-chord-line factor and the fin--body interference factor
  $K_{fb} = 1 + r_b/(s + r_b)$ \src{model/parts/FinSet.cpp}{60} (\cref{sec:aero}).
\item[\code{HollowSphere}] --- a uniform-density thick-walled hollow sphere. The constructor
  rejects $r_i \ge r_o$ and $\rho \le 0$ up front, because $r_i = r_o$ would divide by zero
  inside the tensor formula \src{model/parts/HollowSphere.cpp}{33}. Its length is pole to
  pole, $2r_o$ \src{model/parts/HollowSphere.h}{51}; its silhouette is the circle about a
  center at $z = -r_o$, with the cavity as an inner band
  \src{model/parts/HollowSphere.cpp}{53}; and it overrides \code{isSolid()} because the bore
  vanishes at the poles, where fore-plane inference would call it solid
  \src{model/parts/HollowSphere.h}{59}. It doubles as the boot placeholder body a
  \code{RocketModel} constructs before any design is authored or loaded
  \src{model/RocketModel.cpp}{30}.
\item[\code{Motor}] --- a zero-length leaf that wraps a \code{MotorModel} \emph{by value}
  \src{model/parts/Motor.h}{18}, so a burning motor's mass participates in composite mass,
  CM, and inertia like any other part. Its single physics contribution is the
  \code{getMass($t$)} override, which delegates to the wrapped model: loaded weight before
  ignition, falling to casing mass over the burn, constant after burnout
  \src{model/parts/Motor.h}{40} (\cref{sec:motors}). Swapping motors goes through
  \code{setMotorModel()}, which mutates in place --- re-seeding the tensor and mass so the
  tree recomputes --- rather than re-attaching, so the \code{RocketModel}'s borrowed
  \code{Motor*} handle stays valid \src{model/parts/Motor.cpp}{16}. It is not
  factory-constructible; \code{RocketModel::setMotorModel()} attaches it at the airframe
  root, and since a Motor carries no geometric envelope it never trips the overlap gate
  \src{model/RocketModel.cpp}{118}.
\end{description}

\begin{designrule}[Fin sets are one leaf, not $N$ children]
``The Part tree shares one body-frame orientation and never rotates child tensors (Part.h), so
an azimuthally-arrayed component can't be N children at different angles --- each fin's tensor
would need an Rz rotation the tree won't apply. A FinSet instead bakes the N-fin
rotate-and-sum into its own per-unit-mass tensor'' \src{model/parts/FinSet.h}{25}. This is the
composite walk's omitted $R\,I\,R^{\mathsf{T}}$ rotation (\cref{sec:parts-massprops}) surfacing
as a modeling constraint: the one component whose sub-parts genuinely sit at different
azimuths must pre-sum them analytically in \code{InertiaTensors::TrapezoidalFinSet}
(\cref{sec:parts-tensors}) and present as a single transversely isotropic leaf.
\end{designrule}

\begin{keyinsight}[The motor is the tree's one millimeter island]
Every part API in \srcf{model/parts/} speaks SI meters --- the factory contract says so
explicitly: ``Lengths SI meters, densities kg/m\textasciicircum{}3 --- never mm''
\src{model/parts/PartFactory.h}{25}. The lone exception is the motor's \emph{catalog}
geometry: \code{MotorModel} stores diameter and length in millimeters, the convention of the
RockSim files and the motor database that feed it (\cref{sec:motors}). The unit boundary is
crossed at exactly two lines, both inside \code{Motor::motorTensor()}:
\cppinline|motor.data.diameter / 1000.0| and its length twin, annotated ``RSE / DB store mm''
\src{model/parts/Motor.cpp}{27}. From those two divisions onward everything the tree sees is
meters: the per-unit-mass tensor is a solid cylinder built from the converted dimensions,
degrading to a zero tensor --- a point mass that still contributes its time-varying mass ---
when the catalog geometry is unknown \src{model/parts/Motor.cpp}{31}.
\end{keyinsight}

\subsection{Mass properties: the two-pass composite walk}\label{sec:parts-massprops}

The tree's tensor bookkeeping rests on one convention: a part's \emph{own} stored tensor
(\code{getI()}/\code{setI()}) is
\emph{per unit mass} --- the inertia integral divided by mass, a purely geometric quantity in
$\mathrm{m}^2$ --- while every \emph{composite} tensor (\code{getCompositeI()}) is the full
mass-weighted tensor in $\mathrm{kg\,m}^2$ about the composite CM
\src{model/parts/Part.h}{30}. The constructor seeds the childless composite as $m \cdot I$
\src{model/parts/Part.cpp}{41}, and the composite walk multiplies by \code{getMass($t$)} at
read time \src{model/parts/Part.cpp}{225}. The payoff is that density and geometry are
factored: the closed forms in \code{InertiaTensors} are pure geometry, and a \code{Motor}'s
tensor stays constant while its mass burns down --- the time variation enters once, through
the mass factor. A white-box test pins the convention so it cannot silently regress: for a
2\,kg unit sphere, \code{getI()} reads $0.4$ and \code{getCompositeI()} reads $0.8$
\src{model/tests/PartTests.cpp}{145}.

\begin{physbox}[The parallel-axis (Huygens--Steiner) theorem]
The inertia tensor $I$ measures a rigid body's resistance to angular acceleration; unlike
mass, it depends on the reference point it is taken about, and it is smallest about the
body's own CM. To combine several bodies into one composite, each body's centroidal tensor
must first be re-expressed about the \emph{common} point --- the composite CM. For a body of
mass $m$ whose CM is displaced by the vector $d$ from that point, the theorem gives
\begin{equation*}
I' \;=\; I_{\mathrm{cm}} \;+\; m\big((d \cdot d)\,I_3 - d\,d^{\mathsf{T}}\big),
\end{equation*}
where $I_3$ is the $3\times 3$ identity. The first added term, $m(d \cdot d)I_3$, is the
familiar scalar rule $I = I_{\mathrm{cm}} + m d^2$ applied to every axis at once; the
subtracted $m\,d\,d^{\mathsf{T}}$ removes the component along $d$ itself, because displacing a
body \emph{along} an axis does not change its inertia about that axis. The correction is even
in $d$ --- only the offset matters, not its sign. One subtlety governs composite algorithms:
the map is valid only from a body's own CM to the new point --- shifting through an
intermediate point does \emph{not} add up, so each body must be shifted in one step, from its
own CM directly to the composite CM.
\end{physbox}

All three composite quantities come from one time-aware walk,
\code{computeCompositeAt($t$)} \src{model/parts/Part.h}{230}. It runs over
\code{resolvedCache} --- the resolver output for this sub-tree with the sub-tree root planted
at the local origin \src{model/parts/Part.cpp}{166} --- so geometry is resolved once per
structural change and the walk merely \emph{re-weights} it by \code{getMass($t$)}
\src{model/parts/Part.cpp}{177}. Before any arithmetic, the walk consults the cached placement
verdict and refuses a failed solve (\cref{sec:parts-gating}).

\paragraph{Pass 1: the mass-weighted composite CM.} For each resolved placement, the part's
CM is located in the sub-tree-root frame by composing its resolved fore-plane origin with its
local CM station \src{model/parts/Part.cpp}{196}. Writing the code's own names, with
$\mathtt{pm}_i$ the part's \code{getMass($t$)} and $\mathtt{cmOffset}_i$ the CM-versus-middle
offset of \cref{sec:parts-catalog}:
\begin{align}
\mathtt{cmLocal}_i &= \Big(\mathtt{cmOffset}_{i,x},\; \mathtt{cmOffset}_{i,y},\;
  -\tfrac{L_i}{2} + \mathtt{cmOffset}_{i,z}\Big), \\
\mathtt{cmInRoot}_i &= \mathtt{pose.origin}_i + \mathtt{pose.orient}_i \cdot \mathtt{cmLocal}_i, \\
m = \sum_i \mathtt{pm}_i, \qquad
\mathtt{temp\_cm} &= \frac{1}{m} \sum_i \mathtt{pm}_i \, \mathtt{cmInRoot}_i \qquad (m > 0).
\end{align}
The $-L_i/2$ appears because every part reports its CM relative to the component
\emph{middle}, while its resolved origin is its \emph{fore plane}; the uniform expression
$-L/2 + \mathtt{cmOffset}_z$ converts one to the other for every type with no per-part adapter
\src{model/parts/Part.h}{81}. A fully massless sub-tree has no meaningful CM, so the divide is
guarded and \code{temp\_cm} stays at the origin \src{model/parts/Part.cpp}{210}. Because the
sub-tree root's own pose is the identity, the result is expressed in the root's fore-plane
frame --- for the whole rocket, the nose tip. That \emph{tip datum} is shared by the composite
center-of-pressure fold (\cref{sec:aero}), which is what later makes the static margin a plain
subtraction.

\begin{keyinsight}[Read the composite-CM datum from the tests, not the header]
Three doc comments in \srcf{model/parts/Part.h} still describe the composite CM as ``relative
to this part's own CM'' --- on the \code{CompositeProperties.cm} field
\src{model/parts/Part.h}{90}, on \code{getCompositeCm()} \src{model/parts/Part.h}{104}, and on
the \code{compositeCm} member \src{model/parts/Part.h}{258} --- and \src{sim/Aero.h}{24}
carries the same stale wording. The implementation reports the tip datum: the walk plants the
sub-tree root at the local origin \src{model/parts/Part.cpp}{166} and folds each part at
$\mathtt{pose.origin} + \mathtt{orient}\cdot(-L/2 + \mathtt{cmOffset})$
\src{model/parts/Part.cpp}{203}, so even a childless nonzero-length part reports
$-L/2 + \mathtt{cmOffset}_z$, not zero. The tests state the truth explicitly ---
``Composite CG is now reported in the root's fore-plane (tip) datum, not the root's own CM''
\src{model/tests/BodyTubeTests.cpp}{92} --- and the migration-invariance suite is built around
it \src{tests/PlacementInvarianceTests.cpp}{113}. This is a stale comment, not a code defect
--- the tensor is taken about the composite CM regardless of which datum that CM is reported
in, so no arithmetic is wrong --- but a reader who trusts those three lines will misplace
every CG by half the root's length.
\end{keyinsight}

\paragraph{Pass 2: inertia about the composite CM.} The walk then shifts each part's
mass-weighted tensor to \code{temp\_cm} with the parallel-axis theorem, exactly as the physics
box states it (\cref{lst:parts-passtwo}). The displacement tensor is a three-line helper,
\cppinline|return d.dot(d) * Matrix3::Identity() - d * d.transpose();|
\src{model/parts/Part.cpp}{21}, noted in the source as even in $d$ ``so the sign of d is
irrelevant'' \src{model/parts/Part.cpp}{19}. In the code's names,
\begin{equation}
I \;=\; \sum_i \Big[\, \mathtt{c.mass}_i \cdot I_i \;+\;
   \mathtt{c.mass}_i \big( (d_i \cdot d_i)\, I_3 - d_i\, d_i^{\mathsf{T}} \big) \Big],
\qquad d_i = \mathtt{cmInRoot}_i - \mathtt{temp\_cm},
\label{eq:parts-passtwo}
\end{equation}
where $I_i$ is \code{c.part->getI()}, the per-unit-mass centroidal tensor, so
$\mathtt{c.mass}_i \cdot I_i$ is the part's full centroidal tensor in $\mathrm{kg\,m}^2$ and
the second term is its Huygens--Steiner shift to the composite CM
\src{model/parts/Part.cpp}{225}. The physics box's subtlety --- shift each body from its
\emph{own} CM directly to the composite CM, never through an intermediate point --- is exactly
the bug class the tests pin: a two-point-mass fixture whose inertia about the composite CM is
the reduced-mass value $\mu L^2$ (``the pre-fix code computed this about the PARENT's CM''
\src{model/tests/PartTests.cpp}{199}), a depth-2 chain checked against its flattened
equivalent \src{model/tests/PartTests.cpp}{219}, and the sharpest oracle of all: two coaxial
tubes stacked end-to-end must reproduce one longer tube's mass, CM, and full tensor to
$10^{-12}$ --- the transverse moment grows as $L^2$, so any shift error is loud
\src{model/tests/PartTests.cpp}{308}.

\begin{listing}[htbp]
\begin{minted}{cpp}
   // Pass 2: inertia about the composite CM. Shift each part's own mass-weighted tensor to temp_cm via
   // the parallel-axis theorem. 6-DOF would first rotate the tensor by pl.pose.orient (the R*I*R^T term,
   // identity today, omitted to stay bit-stable).
   // TODO(6-DOF): when that rotation is added, add a test pinning that it reduces to identity in 3-DOF.
   Matrix3 I = Matrix3::Zero();
   for(const Contribution& c : parts)
   {
      const Vector3 d = c.cmInRoot - temp_cm;
      I += c.mass * c.part->getI() + c.mass * parallelAxisTerm(d);
   }
\end{minted}
\caption{Pass 2 of the composite walk \srccap{model/parts/Part.cpp}{217}: the parallel-axis
accumulation of \cref{eq:parts-passtwo}, with the in-tree comment naming the one deliberately
omitted term --- the child-tensor rotation that six-degree-of-freedom flight will
require.}\label{lst:parts-passtwo}
\end{listing}

\begin{keyinsight}[The omitted $R\,I\,R^{\mathsf{T}}$ rotation]
Full generality would rotate each child's tensor into the composite frame before shifting it:
$I_i' = R_i\, I_i\, R_i^{\mathsf{T}}$ with $R_i$ from $\mathtt{pl.pose.orient}$. The code
deliberately omits that step. In three-degree-of-freedom (3-DOF) flight every resolved
orientation is the identity --- the placement layer carries the quaternions but never sets
them (\cref{sec:placement}) --- so the rotation would be an exact no-op, and it is skipped
``to stay bit-stable'' \src{model/parts/Part.cpp}{217}. The omission is total: the only use of
\code{pl.pose.orient} anywhere in the composite path is the CM translation of pass~1
\src{model/parts/Part.cpp}{205}. The class doc calls the inertia shift ``the one piece still
assuming identity'' \src{model/parts/Part.h}{39}, and the companion design document specifies
the one-line 6-DOF form \src{docs/PartPlacementDesignLatex/sections/09-sixdof.tex}{83}. Today's
visible consequence is
structural, not numerical: azimuthally arrayed components cannot be modeled as rotated
children, which is precisely why \code{FinSet} is one analytic leaf
(\cref{sec:parts-catalog}). \Cref{sec:incomplete} carries the consolidated 6-DOF inventory.
\end{keyinsight}

\subsection{Time-awareness and caching}\label{sec:parts-caching}

The single override point for time-varying mass is \code{getMass($t$)}
\src{model/parts/Part.h}{94}. Exactly one production type overrides it --- \code{Motor}
\src{model/parts/Motor.h}{40} --- and everything downstream follows from the composite walk
with no motor-specific inertia code: as propellant burns, the composite mass falls, the
composite CM migrates away from the emptying motor, and the composite tensor re-weights, all
through the same two passes. During a burn the composite tensor tracks an independent
parallel-axis oracle \src{model/tests/MotorTests.cpp}{162}, and after burnout the composites
freeze bitwise \src{model/tests/MotorTests.cpp}{231}. The consumer cadence is per recorded
step: \code{RocketModel::writeMassProperties()} snapshots the triple --- mass, CG,
inertia, ``mutually consistent at this t'' --- into the recorded \code{StateData}
\src{model/RocketModel.cpp}{53}, the seam the 6-DOF rotational ODE will read
\src{sim/StateData.h}{59}; the loop that drives it is \cref{sec:simcore}'s subject.

Freshness is governed by two independent gates with two cadences. The first is structural:
\code{markAsNeedsRecomputing()} sets a per-node flag and walks the parent chain
\src{model/parts/Part.h}{238}, invoked by \code{setMass()}/\code{setI()}
\src{model/parts/Part.h}{74} and by both tree mutators \src{model/parts/Part.cpp}{100}. A
white-box fixture (a declared test-only \code{friend} \src{model/parts/Part.h}{44}) pins the
propagation: a grandchild's \code{setMass()} dirties the root
\src{model/tests/PartTests.cpp}{538}, and removal dirties the ex-parent's whole ancestor chain
\src{model/tests/PartTests.cpp}{585}. The second gate is the mass-delta gate
(\cref{lst:parts-massgate}): absent structural change, the cached CM and tensor rebuild only
when the composite mass has \emph{moved} since the last build.

\begin{listing}[htbp]
\begin{minted}{cpp}
   const double mNow = getCompositeMass(t);
   if(needsRecomputing || mNow != builtAtCompositeMass)
   {
      const CompositeProperties c = computeCompositeAt(t);
      compositeMass          = c.mass;
      compositeCm            = c.cm;
      compositeInertiaTensor = c.inertia;
      builtAtCompositeMass   = mNow;
      needsRecomputing       = false;
   }
\end{minted}
\caption{The mass-delta gate in \code{ensureCompositeCache()}
\srccap{model/parts/Part.cpp}{130}: a structural edit always rebuilds; otherwise the walk
re-runs exactly when the composite mass differs --- by exact floating-point \code{!=} ---
from the mass it was last built at.}\label{lst:parts-massgate}
\end{listing}

\begin{keyinsight}[Why exact \code{==} on a double is the right gate]
The gate key \code{builtAtCompositeMass} is seeded with a quiet NaN, and NaN compares unequal
to everything --- so the first read always builds \src{model/parts/Part.h}{248}. During a
burn, \code{getMass($t$)} differs at every query, so the CM and tensor re-weight every step.
After burnout, the motor's mass function returns the \emph{bit-identical} empty mass for every
$t$, so \code{mNow == builtAtCompositeMass} holds exactly and the cache freezes, bit-stably,
with zero recomputation for the rest of the flight \src{model/parts/Part.cpp}{132}. Keying on
mass rather than on time is deliberate: an adaptive integrator evaluates repeated and
\emph{rejected} stage times (\cref{sec:simcore}), and a time-keyed cache would thrash on them,
while a mass-keyed one is immune \src{model/parts/Part.h}{108}. Two couplings are worth
stating plainly: the gate keys on \emph{total} mass alone, so a hypothetical second
time-varying part whose burn kept the total constant while shifting the CM would be served
stale values --- unrealizable under the single-motor assumption
(\src{model/RocketModel.cpp}{13}; \cref{sec:motors}), but coupled to it; and the freeze
relies on post-burnout \code{getMass} being idempotent float math --- a future mass model
without that property would silently degrade the cache to rebuild-every-step, a performance
hazard rather than a correctness one.
\end{keyinsight}

The gate's key, \code{getCompositeMass($t$)}, is itself never cached: it is a cheap live
recursive sum --- this node plus every descendant, no tensor work
\src{model/parts/Part.cpp}{118} --- because it is both the gate key and the hot divisor in
the equation of motion. The placement cache, finally, runs on
the structural flag alone: \code{placementDirty} is set by add/remove, never by a mass or
inertia edit --- geometry is invariant under a burn \src{model/parts/Part.h}{264}. The two
cadences are witnessed directly: over a simulated burn, an instrumented tube's
radius-query counter stays frozen while the composite CM and mass move every step
\src{model/tests/DiagnosticsGateTests.cpp}{97}. \Cref{sec:placement} covers what that
placement cache holds.

\subsection{The closed-form tensor library}\label{sec:parts-tensors}

The concrete parts take their per-unit-mass tensors from \code{model::InertiaTensors}, a
static-method class of closed forms, each diagonal, about the body's own CM, with $z$ the
symmetry axis \src{model/InertiaTensors.h}{11} (\cref{tab:parts-tensors}).

\begin{table}[htbp]
\centering\small
\begin{tabularx}{\linewidth}{@{}>{\raggedright\arraybackslash}Xlll@{}}
\toprule
Shape & $I_{xx}/m = I_{yy}/m$ & $I_{zz}/m$ & \\
\midrule
Solid sphere, radius $r$ & $\tfrac{2}{5}r^2$ & $\tfrac{2}{5}r^2$ & \src{model/InertiaTensors.h}{17} \\
\addlinespace
Thin-shell sphere, radius $r$ & $\tfrac{2}{3}r^2$ & $\tfrac{2}{3}r^2$ & \src{model/InertiaTensors.h}{28} \\
\addlinespace
Thick hollow sphere, $r_i$, $r_o$ &
  $\tfrac{2}{5}\dfrac{r_o^5 - r_i^5}{r_o^3 - r_i^3}$ &
  $\tfrac{2}{5}\dfrac{r_o^5 - r_i^5}{r_o^3 - r_i^3}$ & \src{model/InertiaTensors.h}{45} \\
\addlinespace
Tube (solid cylinder at $r_i{=}0$) &
  $\tfrac{1}{12}\big(3(r_i^2{+}r_o^2) + L^2\big)$ &
  $\tfrac{1}{2}(r_i^2{+}r_o^2)$ & \src{model/InertiaTensors.h}{63} \\
\addlinespace
Solid cone (CM at $L/4$ from base) &
  $\tfrac{3}{20}R^2 + \tfrac{3}{80}L^2$ & $\tfrac{3}{10}R^2$ & \src{model/InertiaTensors.h}{80} \\
\addlinespace
Conical shell (CM at $L/3$ from base) &
  $\tfrac{1}{4}R^2 + \tfrac{1}{18}L^2$ & $\tfrac{1}{2}R^2$ & \src{model/InertiaTensors.h}{97} \\
\addlinespace
Trapezoidal fin set & \cref{eq:parts-finset} & \cref{eq:parts-finset} & \src{model/InertiaTensors.cpp}{11} \\
\bottomrule
\end{tabularx}
\caption{The closed-form tensor library: per-unit-mass ($\mathrm{m}^2$), diagonal, about each
body's own CM, $z$ the symmetry axis. The thick-sphere form reduces to the solid sphere as
$r_i \to 0$ and to the thin shell as $r_i \to r_o$ \srccap{model/InertiaTensors.h}{39}; the
conical shell's $\tfrac{1}{18}L^2$ term is noted in the source as verified
exact.}\label{tab:parts-tensors}
\end{table}

The one non-textbook entry is the trapezoidal fin set. The model treats one fin as a uniform
trapezoidal lamina in a meridian plane, thickness circumferential, and forms the set tensor as
the azimuthal average of one fin's tensor over the $N$ mounting angles $2\pi k/N$; for
$N \ge 3$ that average is transversely isotropic \emph{and independent of $N$} --- which is
why the fin count parameter is \code{[[maybe\_unused]]} in the implementation
\src{model/InertiaTensors.cpp}{11}. With $\mathit{sum} = c_r + c_t$, the code computes the
radial mass centroid of one fin, $y_c = \tfrac{s}{3}(c_r + 2c_t)/\mathit{sum}$, and its lever
to the body axis, $d = r_b + y_c$ \src{model/InertiaTensors.cpp}{26}, then the single-fin
centroidal second moments --- spanwise $P$ (sweep-independent), chordwise $Q$
(sweep-dependent through the LE sweep $X_t$, as the mean square minus the squared centroid),
and through-thickness $T$ ---
\begin{gather*}
P = \frac{s^2\,(c_r^2 + 4c_r c_t + c_t^2)}{18\,\mathit{sum}^2}, \qquad
x_c = \frac{X_t(c_r + 2c_t) + c_r^2 + c_r c_t + c_t^2}{3\,\mathit{sum}}, \\
\mathtt{meanXsq} = \frac{X_t^2(c_r + 3c_t) + X_t(c_r^2 + 2c_r c_t + 3c_t^2)
   + \mathit{sum}\,(c_r^2 + c_t^2)}{6\,\mathit{sum}}, \qquad
Q = \mathtt{meanXsq} - x_c^2, \qquad T = \frac{t^2}{12}
\end{gather*}
\src{model/InertiaTensors.cpp}{29}, and finally the azimuthal average, in which the
radial--axial product and $z$-coupling terms cancel for $N \ge 3$:
\begin{equation}
I_{zz}/m = P + T + d^2, \qquad
I_{xx}/m = I_{yy}/m = \tfrac{1}{2}\big(P + 2Q + T + d^2\big)
\qquad \text{\src{model/InertiaTensors.cpp}{45}.}
\label{eq:parts-finset}
\end{equation}

The library's acceptance gate is not the algebra but a set of independent numeric oracles ---
the implementation says so itself: ``The mesh oracle in InertiaTensorsTests, not this closed
form, is the acceptance gate'' \src{model/InertiaTensors.cpp}{15}. The solid cone is checked
against a 400{,}000-slice disk-stack integral, with the literal constants $\tfrac{3}{10}R^2$
and $\tfrac{3}{20}R^2 + \tfrac{3}{80}L^2$ pinned at $10^{-15}$
\src{model/tests/InertiaTensorsTests.cpp}{147}; the conical shell against a surface integral
\src{model/tests/InertiaTensorsTests.cpp}{169}; and the fin set against a volume-mesh
integration at a 0.5\% tolerance chosen between the mesh's own error ($\sim$0.01\%) and the
1.2--11\% signature of a wrong (rectangle-only) chordwise term
\src{model/tests/InertiaTensorsTests.cpp}{186}. Cross-checks tighten the net: $I_{zz}$ must
equal an independently derived closed form $J + t^2/12$ at $10^{-12}$
\src{model/tests/InertiaTensorsTests.cpp}{139}, the tensor must be identical across
$N \in \{3, 4, 6\}$ \src{model/tests/InertiaTensorsTests.cpp}{227}, and the rectangle
degeneracy ($c_t = c_r$, $X_t = 0$) must recover $P = s^2/12$, $Q = c^2/12$
\src{model/tests/FinSetTests.cpp}{139}. The mesh also demonstrates the approximation's honest
edge: a 2-fin set is genuinely anisotropic ($I_{xx} \ne I_{yy}$ by more than 1\%), while the
helper knowingly returns the forced-isotropic $N \ge 3$ tensor --- the case the \code{FinSet}
constructor warns about \src{model/tests/InertiaTensorsTests.cpp}{259}.

\subsection{Diagnostics gating on the composite readers}\label{sec:parts-gating}

The composite walk opens with a gate, not a computation. \code{computeCompositeAt()} first
ensures the placement cache is fresh, then reads the cached overlap verdict; if the sub-tree's
geometry failed its solve, it throws a \code{std::runtime\_error} carrying the first located
diagnostic --- ``A self-intersecting design is a hard, located failure, never a
silently-wrong mass/inertia'' \src{model/parts/Part.cpp}{182}. The check costs a flag read,
not a re-sweep: the verdict is computed once per structural change and cached alongside the
resolved geometry \src{model/parts/Part.cpp}{160}.

\begin{designrule}[Fail closed --- except for the ODE divisor]
The gate is asymmetric by design. \code{getCompositeCm()} and \code{getCompositeI()} both pass
through the gated cache and therefore throw on a failed solve, and the test suite pins both
refusals \src{model/tests/DiagnosticsGateTests.cpp}{78}. \code{getCompositeMass()} deliberately
does \emph{not}: it is the hot divisor in the equation of motion, queried at every integrator
stage, and it performs no placement-dependent work --- a mass sum is correct whether or not
two parts interpenetrate. The same tests pin that it keeps answering on a failed design
\src{model/tests/DiagnosticsGateTests.cpp}{83}. The verdict itself is exposed read-only
through \code{placementDiagnostics()} \src{model/parts/Part.h}{125}, and the visualizer reads
that same cached object to paint the offender \src{visualizer/RocketMesh.cpp}{440} --- one
verdict, two consumers, so the simulator's refusal and the on-screen highlight can never
disagree \src{model/parts/Part.h}{274}. What the sweep actually checks --- intervals,
breakpoints, the solid-host rule --- is \cref{sec:placement}'s subject.
\end{designrule}

One more composite fold rides on the same resolved geometry, and it is worth naming here
because it completes the tree's API. \code{getCompositeAero(refArea)} folds each part's
Barrowman contribution additively over the resolved placements, re-expressing every part's
center-of-pressure moment onto the tip datum --- the same datum as the composite CM, so that
the static margin is literally \code{cp() $-$ cg()} \src{model/parts/Part.cpp}{231}. It is a
separate, pure-geometry pass from the mass walk --- it reads no time-varying state
\src{model/parts/Part.h}{119}. At the pinned commit it is also dormant: built, correct on its
tests, and consumed by no production code path --- the test file says so in as many words
(``nothing in production consumes this in P2'' \src{model/tests/AeroTests.cpp}{18}). The
mathematics it folds, the datum trick that makes the fold additive, and what consuming it
would take are the subject of \cref{sec:aero}; \cref{sec:incomplete} places it in the
consolidated inventory.
